\section*{Abstract}
\addcontentsline{toc}{section}{Abstract}
This report presents the design and implementation of a real-time Internet of
Things system on the ESP32-S3 platform using FreeRTOS and PlatformIO. The
system was developed from an RTOS-based starter project and extended into a
multi-functional environmental monitoring node that combines local sensing,
real-time visual feedback, a web-based user interface, on-device TinyML
prediction, and CoreIOT cloud connectivity. The final implementation is not a
single-purpose sensor demo; instead, it is a coordinated embedded system in
which independent tasks exchange data through queues, semaphores, and mutexes.

The implemented node periodically acquires temperature and humidity from a
DHT20 sensor, drives a discrete LED according to temperature conditions,
updates a NeoPixel according to humidity level, and renders context-aware
information on a 16x2 I\textsuperscript{2}C LCD. In parallel, a web server
supports both Access Point (AP) and Station (STA) operation. In AP mode, the
board exposes a dashboard for live monitoring, Wi-Fi scanning, saved-profile
management, and direct control of an external LED through on/off and
brightness-adjustment actions. In STA mode, the node connects to the Internet
and publishes telemetry to the CoreIOT platform using MQTT. The project also
includes a TinyML task that performs batched inference on normalized
temperature-humidity features and generates a weather-oriented rain/no-rain
prediction.

From a software-engineering perspective, the project emphasizes task
separation, predictable timing, and reduced dependence on globally shared
state. Sensor data are treated as streaming information and distributed through
single-slot queues to the LCD, web, TinyML, and cloud modules. An application
context structure stores shared synchronization objects and a compact
shared-state block protected by a mutex. This approach keeps the overall
architecture modular while preserving the current logic of the system.

The report explains the assignment requirements, maps them to the implemented
features, describes the task architecture and data flow, and analyzes the
timing behavior of each subsystem. Based on the code structure and diagnostic
logs, the implementation satisfies the assignment as a soft real-time IoT
system. It achieves periodic sensing, event-driven user interaction, Wi-Fi
mode switching, lightweight on-device inference, and successful preparation for
cloud telemetry publishing in STA mode.

\clearpage
